% GENERATED FILE. Edit canonical lessons, then run npm run generate:books.
% canonical-source-hash: fnv1a64:9caa1fde7aa1dd5f
% canonical-lessons: KA-C491-koragu, KA-C491-marugu, KA-C491-sankata, KA-C491-visada, KA-C491-pascattapa

\chapter{Grief and Regret}
\label{ch:ka-grief-and-regret}
\hlchaptermodality{491}

\begin{chapteropening}
\textbf{By the end of this chapter:} \emph{I can say to pine, to feel pity, anguish, regret and remorse.}

Complete the last lesson of chapter 491: I can say to pine, to feel pity, anguish, regret and remorse.
\end{chapteropening}

\section[\texorpdfstring{\kn{ಕೊರಗು} (koragu)}{ಕೊರಗು (koragu)}]{\kn{ಕೊರಗು} (koragu) — to pine; to brood}

\label{lesson:KA-C491-koragu}



\begin{quote}
Before the new one: say the Kannada for embarrassment, then the Kannada for unease.
\end{quote}

\subsection*{You'll want to know: \kn{ಕೊರಗು}}
\textbf{\kn{ಕೊರಗು}} — \emph{koragu} — \textquotedblleft{}to pine; to brood\textquotedblright{}.

\begin{quote}
\kn{ಕೆಲಸ} \kn{ಹೋದ} \kn{ಮೇಲೆ} \kn{ಅವನು} \kn{ದಿನವೂ} \kn{ಕೊರಗುತ್ತಿದ್ದನು}.
\end{quote}

\begin{quote}
\emph{kelasa hōda mēle avanu dinavū koraguttiddanu.}
\end{quote}

\begin{quote}
After he lost his job, he used to brood every day.
\end{quote}

Past (I …): \kn{ಕೊರಗಿದೆ}, \emph{koragide}.

A quiet, long sorrow kept inside: \kn{ಮನಸ್ಸಿನಲ್ಲೇ} \kn{ಕೊರಗು}, to grieve in silence.

\begin{cousinweb}
\textbf{\kn{ಕೊರಗು}} (\emph{koragu}) is native Dravidian: DEDR 1778 gives Kannada \emph{koragu}, \textquotedblleft{}to lose its sap, shrivel, waste away\textquotedblright{}, and also \textquotedblleft{}to sorrow\textquotedblright{}. A grief that dries a person up is the picture behind it.
\end{cousinweb}

\begin{grammarlens}[title={What you've built}]
One more word for describing what you lived and felt.
\end{grammarlens}

\subsection*{Your turn}
\begin{itemize}
\raggedright
  \item \emph{Say it:} \emph{koragu}
  \item \emph{Say it:} \emph{koragu}, once more
  \item \emph{Say it:} \emph{kelasa hōda mēle avanu dinavū koraguttiddanu.} — After he lost his job, he used to brood every day.
  \item \emph{Recall:} say the Kannada for embarrassment, then the Kannada for unease, then say \emph{koragu} again
\end{itemize}

\begin{tcolorbox}[breakable,colback=teal!4,colframe=teal!35!black,title={Before you move on}]
What does \kn{ಕೊರಗು} mean? (\textquotedblleft{}To pine; to brood\textquotedblright{}.) Say it once more.
\end{tcolorbox}
\section[\texorpdfstring{\kn{ಮರುಗು} (marugu)}{ಮರುಗು (marugu)}]{\kn{ಮರುಗು} (marugu) — to feel pity; to grieve for}

\label{lesson:KA-C491-marugu}



\begin{quote}
Before the new one: say the Kannada for unease, then the Kannada for to pine.
\end{quote}

\subsection*{You'll want to know: \kn{ಮರುಗು}}
\textbf{\kn{ಮರುಗು}} — \emph{marugu} — \textquotedblleft{}to feel pity; to grieve for\textquotedblright{}.

\begin{quote}
\kn{ಗಾಯಗೊಂಡ} \kn{ಹಕ್ಕಿಯನ್ನು} \kn{ನೋಡಿ} \kn{ಮಗು} \kn{ಮರುಗಿತು}.
\end{quote}

\begin{quote}
\emph{gāyagoṇḍa hakkiyannu nōḍi magu marugitu.}
\end{quote}

\begin{quote}
Seeing the injured bird, the child felt sorry for it.
\end{quote}

Past (I …): \kn{ಮರುಗಿದೆ}, \emph{marugide}.

The one you feel for takes -\kn{ಗಾಗಿ}: \kn{ಅವನಿಗಾಗಿ} \kn{ಮರುಗು}, to feel for him.

\begin{cousinweb}
\textbf{\kn{ಮರುಗು}} (\emph{marugu}) is native Dravidian: DEDR 4769 lists Kannada \emph{maṟugu}, \textquotedblleft{}to burn, be very hot\textquotedblright{}, and from there \textquotedblleft{}to suffer, be distressed, grieve\textquotedblright{}. Pity is pictured as a burning inside.
\end{cousinweb}

\begin{grammarlens}[title={What you've built}]
One more word for describing what you lived and felt.
\end{grammarlens}

\subsection*{Your turn}
\begin{itemize}
\raggedright
  \item \emph{Say it:} \emph{marugu}
  \item \emph{Say it:} \emph{marugu}, once more
  \item \emph{Say it:} \emph{gāyagoṇḍa hakkiyannu nōḍi magu marugitu.} — Seeing the injured bird, the child felt sorry for it.
  \item \emph{Recall:} say the Kannada for unease, then the Kannada for to pine, then say \emph{marugu} again
\end{itemize}

\begin{tcolorbox}[breakable,colback=teal!4,colframe=teal!35!black,title={Before you move on}]
What does \kn{ಮರುಗು} mean? (\textquotedblleft{}To feel pity; to grieve for\textquotedblright{}.) Say it once more.
\end{tcolorbox}
\section[\texorpdfstring{\kn{ಸಂಕಟ} (saṅkaṭa)}{ಸಂಕಟ (saṅkaṭa)}]{\kn{ಸಂಕಟ} (saṅkaṭa) — anguish; distress}

\label{lesson:KA-C491-sankata}



\begin{quote}
Before the new one: say the Kannada for to pine, then the Kannada for to feel pity.
\end{quote}

\subsection*{You'll want to know: \kn{ಸಂಕಟ}}
\textbf{\kn{ಸಂಕಟ}} — \emph{saṅkaṭa} — \textquotedblleft{}anguish; distress\textquotedblright{}.

\begin{quote}
\kn{ಅಜ್ಜಿಯ} \kn{ಕಾಯಿಲೆ} \kn{ನೋಡಿ} \kn{ನನಗೆ} \kn{ಸಂಕಟ} \kn{ಆಯಿತು}.
\end{quote}

\begin{quote}
\emph{ajjiya kāyile nōḍi nanage saṅkaṭa āyitu.}
\end{quote}

\begin{quote}
Seeing grandmother's illness, I felt anguish.
\end{quote}

A tight, pressing pain of the mind, and also a tight spot: \kn{ಸಂಕಟದಲ್ಲಿ} \kn{ಸಿಕ್ಕಿಕೊ}, to be caught in a crisis.

\begin{cousinweb}
\textbf{\kn{ಸಂಕಟ}} (\emph{saṅkaṭa}) is borrowed from Sanskrit \emph{saṃkaṭa}, \textquotedblleft{}brought together, narrow, crowded; dangerous, critical\textquotedblright{} (Monier-Williams). Kannada uses it for a narrow place of the mind: distress.
\end{cousinweb}

\begin{grammarlens}[title={What you've built}]
One more word for describing what you lived and felt.
\end{grammarlens}

\subsection*{Your turn}
\begin{itemize}
\raggedright
  \item \emph{Say it:} \emph{saṅkaṭa}
  \item \emph{Say it:} \emph{saṅkaṭa}, once more
  \item \emph{Say it:} \emph{ajjiya kāyile nōḍi nanage saṅkaṭa āyitu.} — Seeing grandmother's illness, I felt anguish.
  \item \emph{Recall:} say the Kannada for to pine, then the Kannada for to feel pity, then say \emph{saṅkaṭa} again
\end{itemize}

\begin{tcolorbox}[breakable,colback=teal!4,colframe=teal!35!black,title={Before you move on}]
What does \kn{ಸಂಕಟ} mean? (\textquotedblleft{}Anguish; distress\textquotedblright{}.) Say it once more.
\end{tcolorbox}
\section[\texorpdfstring{\kn{ವಿಷಾದ} (viṣāda)}{ವಿಷಾದ (viṣāda)}]{\kn{ವಿಷಾದ} (viṣāda) — regret; sorrow}

\label{lesson:KA-C491-visada}



\begin{quote}
Before the new one: say the Kannada for to feel pity, then the Kannada for anguish.
\end{quote}

\subsection*{You'll want to know: \kn{ವಿಷಾದ}}
\textbf{\kn{ವಿಷಾದ}} — \emph{viṣāda} — \textquotedblleft{}regret; sorrow\textquotedblright{}.

\begin{quote}
\kn{ಆ} \kn{ತಪ್ಪಿಗೆ} \kn{ನನಗೆ} \kn{ತುಂಬಾ} \kn{ವಿಷಾದ} \kn{ಇದೆ}.
\end{quote}

\begin{quote}
\emph{ā tappige nanage tumbā viṣāda ide.}
\end{quote}

\begin{quote}
I deeply regret that mistake.
\end{quote}

Formal and polite: notices say \kn{ವಿಷಾದಿಸುತ್ತೇವೆ}, we regret.

\begin{cousinweb}
\textbf{\kn{ವಿಷಾದ}} (\emph{viṣāda}) is borrowed from Sanskrit \emph{viṣāda}, \textquotedblleft{}a drooping state, languor\textquotedblright{} and so dejection (Monier-Williams); in Kannada it is the sorrow of regret.
\end{cousinweb}

\begin{grammarlens}[title={What you've built}]
One more word for describing what you lived and felt.
\end{grammarlens}

\subsection*{Your turn}
\begin{itemize}
\raggedright
  \item \emph{Say it:} \emph{viṣāda}
  \item \emph{Say it:} \emph{viṣāda}, once more
  \item \emph{Say it:} \emph{ā tappige nanage tumbā viṣāda ide.} — I deeply regret that mistake.
  \item \emph{Recall:} say the Kannada for to feel pity, then the Kannada for anguish, then say \emph{viṣāda} again
\end{itemize}

\begin{tcolorbox}[breakable,colback=teal!4,colframe=teal!35!black,title={Before you move on}]
What does \kn{ವಿಷಾದ} mean? (\textquotedblleft{}Regret; sorrow\textquotedblright{}.) Say it once more.
\end{tcolorbox}
\section[\texorpdfstring{\kn{ಪಶ್ಚಾತ್ತಾಪ} (paścāttāpa)}{ಪಶ್ಚಾತ್ತಾಪ (paścāttāpa)}]{\kn{ಪಶ್ಚಾತ್ತಾಪ} (paścāttāpa) — remorse; repentance}

\label{lesson:KA-C491-pascattapa}



\begin{quote}
Before the new one: say the Kannada for anguish, then the Kannada for regret.
\end{quote}

\subsection*{You'll want to know: \kn{ಪಶ್ಚಾತ್ತಾಪ}}
\textbf{\kn{ಪಶ್ಚಾತ್ತಾಪ}} — \emph{paścāttāpa} — \textquotedblleft{}remorse; repentance\textquotedblright{}.

\begin{quote}
\kn{ಸುಳ್ಳು} \kn{ಹೇಳಿದ್ದಕ್ಕೆ} \kn{ಅವನಿಗೆ} \kn{ಪಶ್ಚಾತ್ತಾಪ} \kn{ಆಯಿತು}.
\end{quote}

\begin{quote}
\emph{suḷḷu hēḷiddakke avanige paścāttāpa āyitu.}
\end{quote}

\begin{quote}
He felt remorse for having lied.
\end{quote}

For what: -\kn{ಕ್ಕೆ}. \kn{ಪಶ್ಚಾತ್ತಾಪ} \kn{ಪಡು}, to repent.

\begin{cousinweb}
\textbf{\kn{ಪಶ್ಚಾತ್ತಾಪ}} (\emph{paścāttāpa}) is borrowed from Sanskrit \emph{paścāttāpa}, literally \textquotedblleft{}after-pain\textquotedblright{}: sorrow, regret, repentance (Monier-Williams). It is built from \emph{paścāt}, \textquotedblleft{}afterwards\textquotedblright{}, and \emph{tāpa}, \textquotedblleft{}heat, pain\textquotedblright{}.
\end{cousinweb}

\begin{grammarlens}[title={What you've built}]
One more word for describing what you lived and felt.
\end{grammarlens}

\subsection*{Your turn}
\begin{itemize}
\raggedright
  \item \emph{Say it:} \emph{paścāttāpa}
  \item \emph{Say it:} \emph{paścāttāpa}, once more
  \item \emph{Say it:} \emph{suḷḷu hēḷiddakke avanige paścāttāpa āyitu.} — He felt remorse for having lied.
  \item \emph{Recall:} say the Kannada for anguish, then the Kannada for regret, then say \emph{paścāttāpa} again
\end{itemize}

\begin{tcolorbox}[breakable,colback=teal!4,colframe=teal!35!black,title={Before you move on}]
What does \kn{ಪಶ್ಚಾತ್ತಾಪ} mean? (\textquotedblleft{}Remorse; repentance\textquotedblright{}.) Say it once more.
\end{tcolorbox}
